\section{The learning rule}
\label{sec:rule}

Marginalising both corruptions of \S\ref{sec:setup} gives the likelihood of the
received opinion, and projecting the resulting posterior back onto the Gaussian
family of Eq.~\ref{eq:state} by maximum entropy gives a closed-form online update.
The derivation is standard and deferred to Appendix~\ref{app:derivation}; what
matters in the main text is the result and two of its properties.

\paragraph{Two fields.} Everything the update depends on enters through two
scalars, each a raw field divided by an \emph{internal response uncertainty}:
\begin{equation}
  \hw = \frac{(\hat{\wv}_r \cdot \xv)\,\sigma_e}{\gamma_C},
  \qquad
  \hmu = \frac{\mu_{e|r}}{\gamma_V},
  \qquad
  \gamma_C = \sqrt{1 + \xv \cdot C_r \xv},
  \quad
  \gamma_V = \sqrt{1 + V_{e|r}}.
  \label{eq:fields}
\end{equation}
$\hw > 0$ means the receiver already agrees with the message; $\hmu > 0$ means it
distrusts the source.

\paragraph{The evidence is an exclusive-or.} The marginal likelihood reduces to
\begin{equation}
  Z \;=\; \Phi(\hw) + \Phi(\hmu) - 2\,\Phi(\hw)\Phi(\hmu),
  \label{eq:evidence}
\end{equation}
which is $P(a) + P(b) - 2P(a)P(b)$: the probability that exactly one of the two
holds. $Z$ is large in the two \emph{consonant} configurations --- agree with a
trusted source, disagree with a distrusted one --- and collapses to zero in the
two \emph{dissonant} ones. Surprise, in this model, is not an assumption; it is
what the marginal likelihood does.

\paragraph{The update.} The four modulation functions are the log-derivatives of
Eq.~\ref{eq:evidence},
\begin{equation}
  F_w = \bigl(1 - 2\Phi(\hmu)\bigr)\frac{g(\hw)}{Z},
  \quad
  F_\mu = \bigl(1 - 2\Phi(\hw)\bigr)\frac{g(\hmu)}{Z},
  \quad
  F_C = -F_w(F_w + \hw),
  \quad
  F_V = -F_\mu(F_\mu + \hmu),
  \label{eq:modulation}
\end{equation}
with $g$ the standard normal density, and the receiver moves
\begin{equation}
  \boxed{
  \begin{aligned}
    \hat{\wv}_r &\mathrel{+}= \frac{F_w}{\gamma_C}\,\sigma_e\, C_r \xv,
      &\qquad C_r &\mathrel{+}= \frac{F_C}{\gamma_C^2}\,(C_r\xv)(C_r\xv)^{\!\top},\\[2pt]
    \mu_{e|r}   &\mathrel{+}= \frac{F_\mu}{\gamma_V}\,V_{e|r},
      &\qquad V_{e|r} &\mathrel{+}= \frac{F_V}{\gamma_V^2}\,V_{e|r}^2 .
  \end{aligned}}
  \label{eq:update}
\end{equation}
The opinion update is Hebbian in $\sigma_e \xv$, with a tensorial, adaptive,
$F_w$-modulated learning rate. Note the symmetry $F_w(x,y) = F_\mu(y,x)$ and
$F_C(x,y) = F_V(y,x)$: the two sectors run the same dynamics with their arguments
exchanged. This is what the scaling by $\gamma_C, \gamma_V$ buys, and it is why
group structure can appear in either sector first (\S\ref{sec:agenda}).

\begin{figure}[t]
  \centering
  \includegraphics[width=0.86\textwidth]{learning_flows}
  \caption{The flow induced on $(\hw, \hmu)$ by one interaction, at
    $D = -d$ (left), $D = 0$ (centre) and $D = +d$ (right). Learning moves the
    receiver out of the dissonant quadrants. At $D = 0$ the two consonant basins
    are symmetric, so the population polarises but the split is uncorrelated with
    the label; a nonzero $D$ tilts the basin boundary, which is the entire
    mechanism of \S\ref{sec:spurious}.}
  \label{fig:flows}
\end{figure}

\subsection{Two properties that do all the work}

\paragraph{Anti-learning.} The prefactor $(1 - 2\Phi(\hmu))$ in $F_w$ is
\emph{negative} whenever the source is distrusted. The receiver then moves its
weights away from what the source transmitted --- it learns the opposite.
Effective antiferromagnetic couplings appear in a system where nothing
antiferromagnetic was put in, and this is the engine of every collective effect we
report. Ablation A2 in \S\ref{sec:ablation} removes exactly this sign and nothing
else.

\paragraph{Blame attribution.} On a surprise, one sector must yield, and
Eq.~\ref{eq:modulation} decides which: the sector whose field is \emph{smaller in
magnitude} absorbs the change, with the crossover at $|\hmu| = |\hw|$. A receiver
with a firm opinion and a weak read on its source revises the source's
reliability; one with a weak opinion and a firm read on its source revises its
opinion. Neither behaviour was specified. Both follow from Eq.~\ref{eq:evidence}.

\paragraph{Annealing.} $F_C$ and $F_V$ are negative almost everywhere, so both
covariances shrink monotonically. Since the covariances are also the step sizes
(\S\ref{sec:setup}), the population does not approach a stationary state so much
as run out of the ability to move: it is annealing toward the frozen
$C, V \to 0$ limit, and every agent becomes progressively harder to change as it
becomes more certain. Two consequences we take seriously rather than treating as
nuisances. \emph{When} the observables are measured is a parameter of the
experiment, and we report it (Appendix~\ref{app:params}) where the source draft
leaves it unspecified. And the \emph{order} in which the two sectors organise is
a real feature of the trajectory rather than an artefact of stopping early, which
is what \S\ref{sec:agenda} measures.

\todo{decide whether the modulation-surface figure (\texttt{modulation\_surfaces},
\texttt{modulation\_contours}) earns main-text space or moves wholesale to
Appendix~\ref{app:derivation}. Current recommendation: appendix. The two
properties above are the content, and Fig.~\ref{fig:flows} already shows them
acting.}
